\chapter{How to Use This Book}

\section*{One idea}
Nobody can study every endgame. Endgame manuals are thick, and most of what they teach you will
rarely need. This book starts from the other end. It asks where strong players actually lose
points in the endgame, and studies those endgames first.

To find out, every endgame in tens of thousands of tournament games between grandmasters and
other strong players was checked by a chess engine. Each time a player turned a won position
into a draw, or a drawn position into a loss, the move was marked as a mistake. Counting those
mistakes shows which endgames are worth your time.

\section*{Part I: where the mistakes are}
Chapter~1 shows which endgames you will meet most often. Chapter~2 shows which ones go wrong
most often once they are on the board. Chapter~3 puts the two together: the endgames at the top
of its list cost the most points, and they set the order of the rest of the book.

\section*{Part II: the endgames}
Each chapter in Part~II takes one endgame from that list. It shows positions in which a strong
player went wrong, from the simplest to the more complex, each followed by the rest of the game.

Use them as exercises. Cover the moves below a diagram and decide what you would play. Then
read on and compare. Every diagram has a link that opens the position on Lichess, where you can
play it out against the computer or check your ideas.

\section*{Why grandmaster mistakes?}
A mistake by a grandmaster is not a beginner's slip. It shows where an endgame is hard even for
people who know it well. If you understand why they went wrong, you will avoid the same mistakes
in your own games.
